%! Unit Launch: Study Design — Unit 1, Lesson 1.0 — Slides (Beamer)
% PILOT SHAPE (2026-09-06), Main Ideas / Notes table (2026-09-12). Converted from the
% group-activity deck 2026-09-14: the Data Detectives frame is gone, the practice-launch frame
% is gone, and the meet-card bib claim is now the debrief's cold check.
% GRADUAL RELEASE deck: title → targets → warm-up → hook (unresolved) → I DO divider → two or
% three frames across the four notes rows (the crux frame flagged redacc) → WE DO (a live
% surface, un-answered) → spoken debrief with the cold check → close & START THE HOMEWORK,
% which is the individual practice. No group activity, no practice launch, no exit ticket.
% There is no spoiler rule: the notes frames ARE the direct instruction.
\documentclass[aspectratio=169, 11pt]{beamer}
\usepackage{saar-beamer}   % provides \ceruleanheader{} and \sectionlabel[color]{}

\newcommand{\redink}[1]{{\color{redacc}\bfseries #1}}

\begin{document}

% TITLE SLIDE (CourseName is not defined in beamer, so the name is literal)
{
\setbeamercolor{background canvas}{bg=cerulean}
\begin{frame}[plain]
  \vspace{2.8cm}\hspace{0.55cm}%
  \begin{minipage}{0.9\textwidth}
    {\color{goldacc}\bfseries\small LESSON 1.0}\\[0.5cm]
    {\color{white}\bfseries\LARGE Unit Launch: Study Design}\\[1.2cm]
    {\color{frostmid}\small Statistical Analysis \& Algebraic Reasoning~$\cdot$~Unit 1}
  \end{minipage}
\end{frame}
}

% ── LEARNING TARGETS ─────────────────────────────────────────────────────────
\begin{frame}[t, plain]
  \ceruleanheader{Today's targets}
  \sectionlabel{Lesson 1.0}
  By the end of today, I can\ldots
  \begin{itemize}\setlength{\itemsep}{5pt}
    \item explain why a number is not \textbf{data} until I know \emph{who} was measured,
          \emph{what} was measured, and in what \emph{units}.
    \item read a data table and name the \textbf{individuals} (rows) and the
          \textbf{variables} (columns).
    \item classify a variable as \textbf{categorical} or \textbf{quantitative} --- and justify
          it from what the column records.
    \item use the \textbf{digit test} on a column of digits, and recognize a
          \textbf{statistical question}.
  \end{itemize}
  \begin{block}{How today works}
    Warm-up $\to$ \textbf{notes together, row by row} $\to$ \textbf{one survey we work as a
    class} $\to$ debrief $\to$ \textbf{you start the homework here, alone}.
  \end{block}
\end{frame}

% ── WARM-UP ──────────────────────────────────────────────────────────────────
\begin{frame}[t, plain]
  \ceruleanheader{Warm-Up --- 5 minutes, silent}
  \sectionlabel{percents, proportions, and the mean}
  \vspace{0.3cm}
  \begin{center}
    \begin{minipage}{0.86\textwidth}
      \begin{enumerate}\setlength{\itemsep}{7pt}
        \item \textbf{84} of Riverbend's \textbf{240} seniors ride the bus. What percent?
        \item A survey of \textbf{50} students found \textbf{18} walk to school. Riverbend has
              \textbf{900} students. Predict how many walk.
        \item Five students spent 12, 15, 9, 20, 14 minutes on homework. Find the mean ---
              then write one sentence saying what your answer describes.
      \end{enumerate}
    \end{minipage}
  \end{center}
  \vspace{0.4cm}
  \begin{center}
    {\color{cerulean}\bfseries 3(b) is the one we read aloud. Your sentence stays on the board
    all period.}
  \end{center}
\end{frame}

% ── HOOK — leave it UNRESOLVED; problem 10 (row 4, The Digit Test) settles it ─
{
\setbeamercolor{background canvas}{bg=cerulean}
\begin{frame}[plain]
  \vspace{1.1cm}
  \begin{center}
    {\color{frostmid}\bfseries\small A SPORTS BLOG ON THE RIVERBEND VARSITY ROSTER}\\[0.7cm]
    {\color{white}\bfseries\large ``The roster's mean jersey number is 14. That is low for a
    varsity squad --- this is a young, inexperienced team.''}\\[0.8cm]
    {\color{white}\small Jerseys: 12, \ 7, \ 23, \ 4, \ 17, \ 21 \qquad
    $\dfrac{84}{6} = 14$ \quad --- the arithmetic is \textbf{correct}.}\\[0.8cm]
    {\color{goldacc}\bfseries\large Is the conclusion correct?}\\[0.5cm]
    {\color{frostmid}\small Vote: yes, no, or can't tell. We settle this at problem 10.}
  \end{center}
\end{frame}
}

% ── I DO — direct instruction begins ─────────────────────────────────────────
{
\setbeamercolor{background canvas}{bg=cerulean}
\begin{frame}[plain]
  \vspace{1.8cm}\hspace{0.8cm}%
  \begin{minipage}{0.86\textwidth}
    {\color{goldacc}\bfseries\small GUIDED NOTES \ \textbullet\ \ I DO}\\[0.7cm]
    {\color{white}\bfseries\Large Open to your Guided Notes page. We work the table row by
    row --- fill each term as we name it, not before.}\\[0.9cm]
    {\color{frostmid}\small Five terms go in the vocabulary box today. Write each one the
    moment we use it.}
  \end{minipage}
\end{frame}
}

% ── ROW 1: DATA IN CONTEXT ───────────────────────────────────────────────────
\begin{frame}[t, plain]
  \ceruleanheader{Data are numbers \emph{with a context}}
  \sectionlabel{notes row 1 $\cdot$ Data in Context $\cdot$ problems 1--2}
  \vspace{4pt}
  \begin{center}
    {\Large\bfseries 12} \qquad {\color{slate}\small tells you nothing.}
  \end{center}
  \vspace{6pt}
  \begin{center}
    {\large ``The \textbf{six players on the Riverbend varsity roster} \ scored a
    \textbf{mean of 12} \ \textbf{points} per player.''}\\[6pt]
    {\small\color{cerulean}\bfseries A = who \qquad B = what \qquad C = units}
  \end{center}
  \vspace{8pt}
  \begin{block}{Problem 2 --- and hold on to it}
    The blog's \textbf{14} has a who (six players) and a what (jersey number). What are its
    \emph{units}? \redink{There are none.} A jersey number is not an amount of anything.
    We are not finished with that yet.
  \end{block}
\end{frame}

% ── ROW 2: INDIVIDUALS AND VARIABLES ─────────────────────────────────────────
\begin{frame}[t, plain]
  \ceruleanheader{Individuals and variables}
  \sectionlabel{notes row 2 $\cdot$ Individual \& Variable $\cdot$ problems 3--4}
  \vspace{2pt}
  \begin{center}
    \small
    \renewcommand{\arraystretch}{1.15}
    \begin{tabular}{@{} l c l c c c @{}}
    \hline
    \textbf{Player} & \textbf{Jersey \#} & \textbf{Position} & \textbf{Points} &
      \textbf{Height (in)} & \textbf{Captain} \\
    \hline
    Nia    & 12 & Guard   & 14 & 64 & Yes \\
    Ruben  &  7 & Forward &  9 & 71 & No  \\
    Sasha  & 23 & Center  &  6 & 75 & No  \\
    Marcus &  4 & Guard   & 19 & 68 & Yes \\
    Talia  & 17 & Forward & 11 & 70 & Yes \\
    Devon  & 21 & Guard   & 13 & 72 & No  \\
    \hline
    \end{tabular}
  \end{center}
  \vspace{6pt}
  \begin{columns}[t]
    \begin{column}{0.47\textwidth}
      \begin{block}{Individual}
        One person or object the data are about --- one \textbf{row}. Here: \redink{6}.
      \end{block}
    \end{column}
    \begin{column}{0.47\textwidth}
      \begin{block}{Variable}
        One characteristic recorded for each --- one \textbf{column}. Here: \redink{5}.
      \end{block}
    \end{column}
  \end{columns}
  \vspace{6pt}
  \begin{center}
    {\small\color{cerulean}\bfseries ``Player'' only tells the rows apart. It is not a
    variable we study.}
  \end{center}
\end{frame}

% ── ROW 3: TWO TYPES — the trap is problem 9, LEFT OPEN ──────────────────────
\begin{frame}[t, plain]
  \ceruleanheader{Two types of variables}
  \sectionlabel{notes row 3 $\cdot$ Two Types $\cdot$ problems 5--9}
  \vspace{4pt}
  \begin{columns}[t]
    \begin{column}{0.47\textwidth}
      \begin{block}{Quantitative --- an amount}
        Measured, with units. Summarize it with a \textbf{mean}.\\[4pt]
        {\small Points: $14, 9, 6, 19, 11, 13 \Rightarrow \dfrac{72}{6} = 12$ points.}
      \end{block}
    \end{column}
    \begin{column}{0.47\textwidth}
      \begin{block}{Categorical --- a label}
        A group name. Summarize it with \textbf{counts and percents}.\\[4pt]
        {\small Captain: 3 of the 6 $\Rightarrow$ 50\% are captains.}
      \end{block}
    \end{column}
  \end{columns}
  \vspace{10pt}
  \begin{center}
    {\large The \redink{type} decides what you are allowed to do with the column.}
  \end{center}
  \vspace{8pt}
  \begin{center}
    {\color{goldacc}\bfseries\large Problem 9 --- vote, do not settle: Jersey \#. Categorical
    or quantitative?}\\[4pt]
    {\small\color{slate} Hands up for each. We write the split on the board and move on.}
  \end{center}
\end{frame}

% ── ROW 4: THE CRUX ──────────────────────────────────────────────────────────
\begin{frame}[t, plain]
  \ceruleanheader{The blog's arithmetic was right}
  \sectionlabel[redacc]{notes row 4 $\cdot$ The Digit Test $\cdot$ problems 10--13 $\cdot$ the whole point of today}
  \vspace{4pt}
  \begin{center}
    {\large $\dfrac{12+7+23+4+17+21}{6} = \dfrac{84}{6} = 14$}\\[8pt]
    {\Large\bfseries\color{redacc} Now find the player who wears 14 on average.}\\[6pt]
    {\small\color{slate} 4 \quad 7 \quad 12 \quad 17 \quad 21 \quad 23 \qquad --- nobody
    stands at 14.}
  \end{center}
  \vspace{8pt}
  \begin{block}{The digit test (PS.DC.1c)}
    The test is \textbf{not} ``does it look like a number?'' The test is:
    \redink{does arithmetic on this variable mean anything?}
    Jersey \#, grade level, ZIP code --- all digits, all \textbf{categorical}.
    \textbf{Problem 10 --- vote again: is the conclusion correct?}
  \end{block}
  \vspace{4pt}
  \begin{center}
    {\small\color{cerulean}\bfseries Problem 13 is the counterweight: Height is digits too
    --- and $420/6 = 70$ inches describes a real amount.}
  \end{center}
\end{frame}

% ── WE DO — a LIVE surface: task and data only, no answers ───────────────────
\begin{frame}[t, plain]
  \ceruleanheader{Guided Practice: the Student Life Survey}
  \sectionlabel[goldacc]{We do $\cdot$ problems 14--17 $\cdot$ you hold the pen}
  \vspace{2pt}
  \begin{center}
    \small
    \renewcommand{\arraystretch}{1.1}
    \begin{tabular}{@{} l c c l c l @{}}
    \hline
    \textbf{Student} & \textbf{Grade} & \textbf{Commute (min)} & \textbf{Lunch} &
      \textbf{Siblings} & \textbf{Club member} \\
    \hline
    Alina   & 10 & 20 & Home       & 2 & Yes \\
    Devon   & 11 &  8 & Cafeteria  & 0 & No  \\
    Marisol &  9 & 35 & Home       & 3 & Yes \\
    Theo    & 12 & 15 & Cafeteria  & 1 & Yes \\
    Priya   & 10 &  5 & Off campus & 2 & No  \\
    Jamal   & 11 & 25 & Cafeteria  & 4 & Yes \\
    \hline
    \end{tabular}
  \end{center}
  \vspace{4pt}
  \begin{itemize}\setlength{\itemsep}{3pt}
    \item[14.] How many individuals? How many variables? Name one of each.
    \item[15.] Classify Commute, Lunch, and Grade. Theo says ``Grade is quantitative ---
          it's numbers.'' Run the test on Grade and fix his reason.
    \item[16.] Find the mean commute --- then say it as data: who, what, units.
    \item[17.] Which one takes data to answer: \emph{How long is Priya's commute?} or
          \emph{How long are Riverbend students' commutes?} Rewrite the other.
  \end{itemize}
\end{frame}

% ── DEBRIEF ──────────────────────────────────────────────────────────────────
\begin{frame}[t, plain]
  \ceruleanheader{Debrief: correct arithmetic, wrong kind of variable}
  \sectionlabel{10 minutes $\cdot$ pens down, papers up --- we say these aloud}
  \vspace{4pt}
  \begin{enumerate}\setlength{\itemsep}{4pt}
    \item The blog's \textbf{14} --- say the sentence: the arithmetic is right, the variable
          is a label, the mean describes nobody.
    \item Problem 15 --- two reasons on the board. Which one also catches Member ID tonight?
    \item Problem 13 --- \emph{Height is digits too. Who does 70 inches describe?}
    \item Problem 17 --- how many answers does \emph{How long is Priya's commute?} have?
  \end{enumerate}
  \vspace{4pt}
  \begin{block}{Cold check --- hands down, thirty seconds}
    At Saturday's cross-country meet five Riverbend runners wore bibs \textbf{31, 18, 45, 9,
    12}. A recap blog wrote: ``Riverbend's mean bib number was \textbf{23} --- a low number
    for a varsity meet. This is a young, inexperienced squad.'' The arithmetic is right.
    \textbf{Is the conclusion?}
  \end{block}
\end{frame}

% ── YOU DO — the homework is the individual practice, started in class ───────
\begin{frame}[t, plain]
  \ceruleanheader{You do: start the homework}
  \sectionlabel[goldacc]{Last 10 minutes $\cdot$ on your own $\cdot$ finish it at home}
  \begin{itemize}\setlength{\itemsep}{5pt}
    \item \textbf{Northgate Recreation Center} --- eight members on a record card.
    \item Count the individuals and the variables; classify four columns \emph{with a reason}.
    \item Mean visits, and the percent taking a group class --- one of those two is the only
          summary its column allows.
    \item A staff member computes the mean \textbf{Member ID} and gets 140.5. What does that
          tell you about Northgate's members?
  \end{itemize}
  \begin{block}{Silent and alone --- I am circulating. Packet pages, scored out of 12, due at the start of the first class after two study halls.}
    Start with \textbf{1, 2, and 6}; 3, 4, and 5 can go home. Item 6 is the one I am reading
    over your shoulder.
  \end{block}
\end{frame}

% ── CLOSE ────────────────────────────────────────────────────────────────────
{
\setbeamercolor{background canvas}{bg=cerulean}
\begin{frame}[plain]
  \vspace{1.5cm}\hspace{0.8cm}%
  \begin{minipage}{0.86\textwidth}
    {\color{goldacc}\bfseries\small BEFORE YOU GO}\\[0.7cm]
    {\color{white}\bfseries\Large You walked in able to compute a mean. You walk out able to
    say when a mean is allowed to exist --- and that is decided by the variable, not by the
    arithmetic.}\\[0.9cm]
    {\color{frostmid}\small Next: Lesson 1.1 --- \emph{The Data Cycle}. The four stages, the
    questions a data set can actually answer, and the frequency table. Today's words are
    assumed from here on.}
  \end{minipage}
\end{frame}
}

\end{document}
